\section{Related Work}
\label{sec:related}

This section reviews seven strands of work that bear on \hammer{}: record linkage, author name disambiguation, benchmarks and ground truth, candidate generation and assignment, evidence combination, selective prediction with human review, and metadata integration with provenance. It is built from a literature inventory compiled for this project. Most entries were checked at the level of bibliographic record and abstract, and a minority were read more closely; statements below that something was not found refer only to that inventory and are not evidence that nothing exists.
The inventory is not a systematic review. It was compiled on 19 November 2026 from the project's earlier audit documents and a small targeted extension of exploratory web searches and bibliographic-record look-ups (Crossref, arXiv, Zenodo, JMLR and PMLR pages); no pre-registered protocol, database export or dual screening was used, and no count of screened records was kept. Scopus, Web of Science, Google Scholar, dblp and Semantic Scholar were not searched. The characterisations of S2AND, WhoIsWho, CONNA, HUMO, r-HUMO and LAGOS-AND below were checked against their abstracts and the readable parts of their text on 20 November 2026; the JASIST version of the LAGOS-AND paper could not be read, so its description rests on the arXiv preprint and the Zenodo record.

\subsection{Record Linkage and Work-Level Matching}
Fellegi and Sunter cast record linkage as a decision among links, non-links and possible links under a probabilistic model of field agreement~\cite{fellegi1969recordlinkage}, and later surveys of duplicate-record detection build on that framing~\cite{elmagarmid2007duplicate}. Work-level deduplication remains an active applied task, for example in the preparation of systematic reviews~\cite{forbes2024deduplicator}. These methods ask whether two \emph{records} describe the same item. They do not decide whether a name refers to a given \emph{person}, which is why \hammer{} keeps the two stages apart. \hammer{}'s fixed-weight score has no estimated match and non-match probabilities and is not calibrated, so it should not be read as a Fellegi--Sunter model, even though its three outcomes resemble the three decision regions. Publication deduplication is part of the system but is not the subject of the research questions studied here.

\subsection{Author Name Disambiguation}
Author name disambiguation (AND) has been surveyed repeatedly~\cite{ferreira2012brief,debonis2023graph}. A recent review of 211 documents groups methods into grouping (similarity and clustering) and assignment (classification) families~\cite{rodrigues2024andreview}. Early learning-based work treated citation-level disambiguation as supervised classification~\cite{han2004twosupervised}, domain-specific methods were developed for MEDLINE~\cite{torvik2009authority}, and probabilistic frameworks were proposed for digital libraries~\cite{tang2012unified}.

More recent systems combine several sources and signals. PubMed Computed Authors provides an open resource of disambiguated PubMed author names~\cite{tian2024pubmedcomputedauthors}. AggAND aggregates several large scholarly databases to improve PubMed author name disambiguation~\cite{zhang2021aggand}. Boukhers and Asundi disambiguate DBLP authors with neural representations of names, coauthors and publication text~\cite{boukhers2024deepdblp}, and Choi et al.\ select a disambiguation scheme according to the metadata available from heterogeneous academic sites~\cite{choi2024heterogeneoussites}. Aggregating sources, using several evidence types, and adapting to available metadata therefore all have prior art, and \hammer{} claims none of them as new.

\subsection{Benchmarks and Ground Truth}
Benchmarks differ in task form and in how identity labels are obtained. S2AND consolidates eight existing datasets into a common format and reports that models trained on a single dataset may generalize poorly to others~\cite{subramanian2021s2and}. WhoIsWho provides a benchmark of more than one million papers with an evaluation task for real-time assignment of new papers to candidate authors~\cite{chen2023whoiswho}. GSD1900 offers 1,900 MEDLINE publication pairs for pairwise evaluation~\cite{vishnyakova2019gsd1900}. LAGOS-AND derives its gold-standard data, including a block subset, from ORCID and DOI evidence~\cite{zhang2023lagosand}, and the version used in this work was released for MAG and OpenAlex author identifiers~\cite{zhang2022lagosandv2}. The quality of AND datasets and the way AND is evaluated have themselves been examined~\cite{muller2017datasets,kim2018dblp}.

These resources are not interchangeable with the task studied here. Most formulate AND as clustering, pairwise classification or paper-to-person assignment, whereas \hammer{} resolves a query, a name with optional context, to an author profile in an external provider's index. Labels derived from ORCID linkage deserve particular care in our setting: the identifier that defines the label may also be present in the candidate source being evaluated, so agreement on it is a provisional matching mechanism and not independent identity truth. We return to this circularity in Sections~\ref{sec:design} and~\ref{sec:threats}.

\subsection{Candidate Generation and Assignment}
Entity-resolution pipelines commonly restrict the pairs to be compared through blocking and filtering, a stage that has its own survey literature~\cite{papadakis2020blocking}. In author disambiguation, assignment of a new paper to existing candidate authors is an explicit task in WhoIsWho~\cite{chen2023whoiswho}. CONNA goes a step further: it matches a paper to its best candidate person and then decides whether to assign the paper or to create a new person, with the matching and decision components fine-tuned jointly by reinforcement learning after separate training~\cite{chen2020conna}. The decision after ranking is therefore not new, and \hammer{}'s gates are best seen as a conventional, rule-based variant.

\hammer{} draws its candidates from OpenAlex~\cite{priem2022openalex}. We treat an OpenAlex author identifier as a provider profile and not as a verified person. Because the pool is whatever the provider returns for a query, it bounds what any later stage can achieve: if the expected profile is absent, ranking and decision rules cannot recover it. Retrieval effectiveness is therefore a separate question (RQ2a). Within the inventory we did not find an AND study that reports a retrieval-stage measure such as Recall@$k$ separately from final assignment quality; the summaries available to us for CONNA and WhoIsWho describe assignment accuracy and do not mention one.

\subsection{Evidence Signals and Their Combination}
Prior systems rely on name, coauthor, publication-text, venue and affiliation evidence in various combinations~\cite{boukhers2024deepdblp,zhang2021aggand,tian2024pubmedcomputedauthors}, and the choice of how to weigh them may be learned or rule-based. In the probabilistic tradition the weights come from estimated agreement and disagreement probabilities~\cite{fellegi1969recordlinkage,tang2012unified}. \hammer{} instead uses fixed engineering weights that are renormalized over the available signals. These weights have not been calibrated or validated, and because candidates may be scored over different signal sets, their scores are not always directly comparable. We state this as a limitation of the method and not as a design virtue (Section~\ref{sec:methodology}).

Institutional evidence deserves a separate note. Resolving raw affiliation strings to organization identifiers is a distinct, tool-supported task~\cite{duran-silva-etal-2024-affilgood,kallipoliti2025affro}. Such tools could supply institutional evidence to an author resolver, but the present work does not evaluate them.

\subsection{Selective Prediction and Human Review}
A predictor that may decline to answer trades coverage for lower error on the answers it gives. This reject-option framing goes back to Chow~\cite{chow1970reject}, was developed theoretically for selective classification~\cite{elyaniv2010foundations}, extended to deep networks with an integrated reject option~\cite{geifman2019selectivenet}, and has been surveyed more recently~\cite{hendrickx2024reject}. Risk and coverage are its standard evaluation quantities, and we adopt that vocabulary for RQ2b. The literature is general, and we did not find it applied to author identity with an explicit manual-review and unresolved split.

Human review in disambiguation is also established. AMiner combines clustering, maintenance and human annotators in the loop~\cite{zhang2018aminerhitl}. For entity resolution more generally, HUMO and r-HUMO divide the workload between machine and human to meet precision and recall targets~\cite{chen2018humo,hou2020rhumo}, and CrowdER uses crowdsourcing to resolve uncertain pairs~\cite{wang2012crowder}. \hammer{}'s manual-review outcome is a routing decision to a human, and the system gives no guarantee on the quality of what it assigns automatically; whether the gates achieve useful precision at useful coverage is an empirical question for RQ2b.

\subsection{Metadata Integration and Provenance}
OpenAlex is an open index of scholarly works, authors, venues, institutions and concepts~\cite{priem2022openalex}. Source coverage differs and depends on the task, and comparisons restricted to works shared by several indexes do not establish global coverage~\cite{culbert2025openalexcoverage}. Combining several databases has been used to improve author disambiguation~\cite{zhang2021aggand,choi2024heterogeneoussites}. The database literature gives provenance a formal treatment in terms of why- and where-provenance~\cite{buneman2001whyandwhere}. \hammer{} records, for each stored publication and candidate, which source supplied it and which evidence supported a decision. We use ``provenance'' in this practical sense and do not claim that the system implements the formal notions.

\subsection{Summary and Positioning}
Table~\ref{tab:literature} summarizes the works discussed. In short, AND is a mature field with several task forms, and multi-source evidence, an assign-or-new decision, human review and risk--coverage evaluation each exist separately. The gap we pursue is narrower than any technique. Within the sources reviewed, it is not established that prior work reports retrieval-stage recall and selective automatic-assignment precision against coverage together, on one frozen candidate pool, with manual-review, unresolved, empty-pool and provider-failure outcomes counted in the denominators. This is a statement about reporting and evaluation design. It is not a claim that \hammer{} is more accurate than any listed system, and no such comparison is made.

% Rows summarise the cited records. S2AND, WhoIsWho, CONNA, HUMO and r-HUMO were checked against abstract and readable text on 20 Nov 2026 (REFERENCE_VERIFICATION.md);
% the other rows rest on abstract-level or bibliographic-record readings. Cells saying "not compared" mean the cited work reports no measure that can be set beside HAMMER's.
\begin{table*}[t]
\caption{Representative prior work and its relation to \hammer{}. Task form, evidence and evaluation are as described by the cited sources.}
\label{tab:literature}
\centering
\footnotesize
\renewcommand{\arraystretch}{1.15}
\begin{tabular}{@{}p{0.15\textwidth}p{0.21\textwidth}p{0.19\textwidth}p{0.16\textwidth}p{0.19\textwidth}@{}}
\toprule
\textbf{Work} & \textbf{Task form} & \textbf{Evidence or ground truth} & \textbf{Evaluation family} & \textbf{Relation to \hammer{}} \\
\midrule
Fellegi and Sunter~\cite{fellegi1969recordlinkage} & Record-pair linkage with link, non-link and possible-link decisions & Probabilistic model over field agreement & Decision error rates & Foundation for the work-level layer; \hammer{}'s score is not a fitted probabilistic model \\
Rodrigues et al.~\cite{rodrigues2024andreview} & Review; grouping versus assignment families & Literature (2003--2022) & Not applicable & Taxonomy used to place \hammer{} \\
S2AND~\cite{subramanian2021s2and} & Clustering of author mentions across eight harmonised datasets & Existing dataset labels & Cluster-level F1 & Not compared: no candidate retrieval or abstention \\
WhoIsWho~\cite{chen2023whoiswho} & Benchmark with a real-time assignment task, more than one million papers & Curated benchmark labels & F1, AUC and MAP, by task & Closest task form; \hammer{} resolves a query to a provider profile, not a paper to a person \\
CONNA~\cite{chen2020conna} & Match a paper to the best existing person, then assign or create a new person & Two datasets; deployed on AMiner & F1 & Precedent for a decision stage after ranking \\
AMiner~\cite{zhang2018aminerhitl} & Clustering, maintenance and human annotators in the loop & Deployed system data & Not compared & Precedent for human review in the loop \\
AggAND~\cite{zhang2021aggand}; PubMed Computed Authors~\cite{tian2024pubmedcomputedauthors} & Disambiguation of PubMed author names with additional databases or multiple signals & PubMed-based datasets & Pairwise accuracy, F1 & Multi-source evidence has prior art; biomedical scope \\
Choi et al.~\cite{choi2024heterogeneoussites} & Scheme chosen according to metadata available from several academic sites & Non-public data & Not compared & Metadata-conditioned evidence has prior art \\
LAGOS-AND~\cite{zhang2023lagosand,zhang2022lagosandv2} & Gold-standard blocks of author records & ORCID and DOI linkage & Dataset, not a method & Benchmark candidate; labels are provisional (Section~\ref{sec:design}) \\
HUMO~\cite{chen2018humo}; r-HUMO~\cite{hou2020rhumo} & Machine--human workload split for entity resolution with quality targets & Entity-resolution datasets & Precision and recall targets against human cost & Quality-controlled review exists; \hammer{} offers no such guarantee \\
Selective classification~\cite{chow1970reject,elyaniv2010foundations,geifman2019selectivenet} & Predictors that may abstain & General classification & Risk--coverage & Supplies the evaluation vocabulary used for RQ2b \\
\bottomrule
\end{tabular}
\end{table*}

